% Validation --- section structure from prd.md 3.1 item 5, ownership from Table 3.
% Scaffold only: headings and planning notes, no prose. Delete each note as its section is filled.
%
% This chapter OWNS: the latency experiment, the keyword-spotter evaluation, the acoustic-
% robustness experiment and the formation-control experiment, each with its statistical analysis
% (Table 3: Exp-2, Exp-3, Exp-4). Every measured figure in this document is defined here; Ch 6
% quotes it exactly.
% It does NOT own: the targets (Ch 1 tables, Ch 3 budget -- quote, never restate differently), the
% speaker-sensitivity WER (Memoire de Master; cite it as an input, one sentence), the model
% comparison (Memoire de Master).
%
% Naming (issue 00): experiments in words (latency experiment, keyword-spotter evaluation,
% acoustic-robustness experiment, formation-control experiment), never Exp-2/3/4; criteria by the
% names in Ch 1's tables, never NFR numbers; reflex/parse path. The roadmap brief (START_SESSION.md
% B10) and results/*.md still use the old names -- translate, do not copy.
%
% Framing already decided (STATE.md, "Exp-2 RESULT" and the Exp-3 entries): targets and verdicts
% are kept as set before measurement; each miss is reported as a miss with its located cause, not
% as a failure of the project and not softened. Discuss whether a target was realistic as
% discussion, never as a changed verdict. The end-to-end re-baseline (ADR / STATE.md) is stated as
% a re-baseline, never as a pass.
%
% Tables: \input the generated tables, never retype them --
%   thesis/generated/exp2_latency_budget.tex (\label{tab:latency-budget}, budget vs measured);
%   thesis/generated/table20_end_to_end.tex (\label{tab:end-to-end}; last column is a STAGE SUM,
%   not end-to-end latency -- say so in the caption or text).
% Before quoting any number, check it against the results/ file named below, not against this note.
%
% Promises made elsewhere that this chapter must keep (issue 06, agents/06-ingenieur-ch1/
% argument.md P1-P31):
%   - the full per-class curve of the keyword spotter with the chosen threshold marked (Ch 1,
%     keyword-spotter paragraph of the safety problem);
%   - the reflex path measured idle AND while the parse path decodes (Reflex latency criterion);
%   - every stage instrumented and reported against Ch 3's per-stage budget (Per-stage
%     attribution);
%   - beside the zero collisions, the number of clamp interventions (Ch 1, collision paragraph);
%   - formation quality measured on its own, not inside the acoustic trials;
%   - the direction of the recognition curve as a separate measurement from the rate;
%   - the per-trial CSV record and the single-command rerun, reported against (Ch 1, the
%     requirements common to both documents);
%   - the safe-failure rate, which Ch 6 returns to;
%   - the WER on the author's voice cited from the Memoire de Master, not restated.

\chapter{Validation}
\label{chap:validation}
This chapter measures the system of Chapter~\ref{chap:implementation} against the requirements of
Section~\ref{sec:engineering-requirements}. Every criterion is taken from
Table~\ref{tab:functional-requirements} or Table~\ref{tab:nonfunctional-requirements}, and every
stage of both paths is judged against the budget of Table~\ref{tab:stage-budget}; each target was
fixed before the measurement that answers it and is quoted here unchanged.
Section~\ref{sec:latency-experiment} reports the latency experiment,
Section~\ref{sec:keyword-spotter-evaluation} the keyword-spotter evaluation,
Section~\ref{sec:acoustic-robustness} the acoustic-robustness experiment and
Section~\ref{sec:formation-control} the formation-control experiment, each with its statistical
analysis. Section~\ref{sec:requirements-summary} collects every verdict in one table. A criterion
is met or missed on the point estimate of the statistic its budget names, with one exception: the
two keyword criteria are rates estimated from few events, and they follow an interval rule declared
before they were scored, stated where they are reported. A missed criterion is reported as missed,
together with the stage or mechanism the measurement locates as its cause.

\paragraph{Device and thermal state.} Every wall-clock figure in this chapter was measured on the
Raspberry~Pi~5, with the \texttt{performance} frequency governor and the threads pinned as in
Table~\ref{tab:core-allocation}. The latency experiment ran with an active cooler fitted and
recorded the processor's temperature and throttle flag at the start of every trial: the temperature
stayed between 50.5 and 77.9~$^{\circ}$C, and the flag was clear on all 360 trials. The
acoustic-robustness experiment ran on the same board and logged the same two quantities every
30~s; its thermal state is given with its figures in Section~\ref{sec:acoustic-robustness}. The
\emph{M\'emoire de Master} measured the same board without active cooling and found it throttled on
every trial, so a timing figure is interpretable only together with the thermal state it was
measured in. The keyword-spotter evaluation and the formation-control experiment report no
wall-clock time: the spotter's detection delay is counted in audio frames, and the formation
trials run in simulated time.

\paragraph{Records and reruns.} Every trial is kept as one record: a CSV row per trial for the
formation-control experiment, and a line of \gls{json} per trial for the latency and
acoustic-robustness experiments, which also carries each segment's transcript, raw model output
and stage timings. Every summary and every table in this chapter is computed from those records
by script; none is typed by hand. The speech recogniser is built from a fixed commit of
\texttt{whisper.cpp}, the model files are identified by their SHA-256 hashes, and the Python
environment on the board is pinned to the workstation's versions. Every random draw --- the noise
excerpt mixed into each recording, the bootstrap resamples, the initial positions of each
formation trial --- comes from a fixed seed, and decoding is greedy. Each experiment is rerun by
one command on the host that produced its records, and one further command regenerates every table
from them. The single-command rerun of Section~\ref{sec:engineering-requirements} therefore holds
per experiment rather than for the chapter as a whole.

\section{Latency experiment}
\label{sec:latency-experiment}
The latency experiment times both paths on the board and judges four criteria of
Table~\ref{tab:nonfunctional-requirements} on the result: reflex latency, end-to-end latency,
per-stage attribution and preemption recovery. It also supplies the preemption trials and the
matrix of detected classes on which the reflex-path requirement of
Table~\ref{tab:functional-requirements} is verified, and the keyword false rejects on the author's
voice. Table~\ref{tab:latency-budget} sets every measured stage beside its target.

\paragraph{Protocol.} Recorded audio is replayed through the runtime on the cooled board, in real
time, through the frame source of Section~\ref{sec:audio-chain}, rather than captured from a
microphone. Replay gives the idle and loaded conditions identical audio, and it makes every trial repeatable. The parse path
is measured on the 200 golden-set utterances in clean audio, one latency sample per endpointed
segment, with the reflex path listening as deployed. The reflex path is measured on the author's 40
recorded keyword takes, 20 per class, each replayed twice in each of two conditions, for 80
trials per condition: idle, with the parse path resident but given no work, and loaded, with a load
thread on cores~1--3 running the deployed parse path back to back on golden-set transcripts, so
that a trigger meets a decode in progress. The loaded condition is preceded by 300~s of the same
load, after which the processor was at 73.0~$^{\circ}$C. The bus runs
over the board's loopback interface, with the consumer and the state machine on the board
(Section~\ref{sec:command-bus}). The reflex-path span ends at publication, as
Table~\ref{tab:stage-budget} defines it. The parse-path span ends later, when the state machine on
the board has applied the command, so it includes the bus and the state-machine check. Both fall
within the validation row of Table~\ref{tab:latency-budget}, whose p95 is under 1~ms, and the
measured span therefore exceeds the defined one by at most that row. The keyword offset is the end
of speech as the capture tool recorded it from the signal's energy, not a forced alignment.

% Generated by eval/tables.py -- do not edit; edit the renderer and re-run.
\begin{table}[htbp]
  \centering
  \footnotesize
  \setlength{\tabcolsep}{4pt}
  \caption[Latency budget versus measured]{Latency budget versus measured performance on Raspberry~Pi~5 across parse and reflex paths, with nearest-rank percentiles and p95 verdict. Bold rows denote non-functional requirement criteria.}
  \label{tab:latency-budget}
  \begin{tabular}{lrrrrrl}
    \toprule
    Stage & n & p50 & p95 & p99 & Target p95 & Verdict \\
    \midrule
    Endpointing wait & 226 & 480 & 480 & 480 & 500 & MEETS \\
    Speech recognition (whisper \texttt{tiny.en}, 3 threads) & 226 & 1,246 & 1,449 & 1,571 & 1,200 & MISSES \\
    Language-model prefill & 226 & 570 & 751 & 849 & 250 & MISSES \\
    Language-model decode & 226 & 561 & 915 & 984 & 1,100 & MEETS \\
    Validation, state-machine check and dispatch & 226 & 0 & 0 & 4 & 50 & MEETS \\
    \textbf{End to end, from end of speech} & 226 & 2,441 & 3,122 & 3,380 & 2,500 & MISSES \\
    End to end, including the endpointing wait & 226 & 2,921 & 3,602 & 3,860 & 2,500 & MISSES \\
    End to end, from start of speech & 226 & 6,344 & 10,086 & 12,322 & -- & -- \\
    \textbf{Reflex path from keyword offset, idle} & 78 & 195 & 545 & 555 & 150 & MISSES \\
    of which spotter delay and frame quantisation & 78 & 190 & 540 & 550 & -- & -- \\
    of which Pi compute and publish & 78 & 5 & 6 & 6 & -- & -- \\
    \textbf{Reflex path from keyword offset, loaded} & 78 & 205 & 547 & 564 & 150 & MISSES \\
    of which spotter delay and frame quantisation & 78 & 190 & 540 & 550 & -- & -- \\
    of which Pi compute and publish & 78 & 12 & 18 & 19 & -- & -- \\
    Reflex path from keyword onset, idle & 78 & 965 & 1,265 & 1,265 & 850 & MISSES \\
    Reflex path from keyword onset, loaded & 78 & 967 & 1,266 & 1,278 & 850 & MISSES \\
    \textbf{Preemption recovery} & 78 & 601 & 1,195 & 1,505 & 300 & MISSES \\
    \bottomrule
  \end{tabular}
  \par\medskip
  \begin{minipage}{\textwidth}\footnotesize\raggedright Prefill is the deployed parser's, with the prompt cache off; the budget's prefill allowance assumes a cached prefix. Start-of-speech has no verdict: its 5,500 ms target is for a 3 s utterance. Split rate, cross-trigger matrix, pre-emption outcomes and the core-0 frame budget: \texttt{results/exp2\_analysis.md}.\end{minipage}
\end{table}


\paragraph{Reflex latency.} The criterion is missed in both conditions: the reflex path's p95 from
the keyword offset is 545~ms idle and 547~ms loaded, against 150~ms. The two sub-rows of
Table~\ref{tab:latency-budget} locate the miss. The spotter's delay and the quantisation of audio
into 80~ms frames account for 540~ms at p95 in both conditions, while the board's computation and
publication account for 6~ms idle and 18~ms loaded. The miss therefore lies in how long the
spotter's score takes to cross the threshold after the keyword ends, a property of the classifier
and its operating point, not of the board. On the synthetic test recordings the same delay is
shorter (Section~\ref{sec:keyword-spotter-evaluation}), so the figure measured here on the author's
voice is the one the criterion is judged on. Loading the parse path raised the reflex path's p95 by
2~ms, and no frame in any condition took longer than the 80~ms frame period to process: in the
loaded condition the frame loop's computation was 14.2~ms at the median, 22.9~ms at p99 and
45.5~ms at most. The core allocation of Section~\ref{sec:resource-allocation} kept the parse path's
load off the reflex path. The onset-anchored figures, 1{,}265~ms idle and 1{,}266~ms loaded against
850~ms, also miss.

\paragraph{Preemption and the matrix of classes.} Of the 80 loaded trials, 78 produced a
detection, and every one of the 78 found a parse-path decode in progress and cancelled it. No
cancelled decode completed despite the cancellation, so the ordering rule of
Section~\ref{sec:command-bus} never had to discard a late result. In both conditions, the 40
\emph{swarm abort} takes all fired the abort class, and 38 of the 40 \emph{swarm hold} takes fired
the hold class. The two misses are one take, missed on both of its replays. No take fired the other
class, in either condition. The preemption behaviour and the class matrix together satisfy the
verification that the reflex-path requirement names. While listening to the 200 golden-set
utterances, the reflex path fired once, a \texttt{hover} during an utterance commanding a landing.
That error is on the fail-safe side of Section~\ref{sec:membership-rule}. The spotter scored the
same recordings without a false accept in the keyword-spotter evaluation, and the records do not
explain the difference.

\paragraph{Preemption recovery.} The criterion is missed: the time from a trigger until the parse
path can accept a new utterance is 1{,}195~ms at p95, against 300~ms. The client's part of the
cancellation is fast, since the cancelled request returns within 6~ms of the trigger at p95. The
remainder is the language-model server's. The server accepts a new request only after it has
registered the closed connection, which it checks once per second. That check, not the
cancellation, sets the recovery time, whose median is 601~ms.

\paragraph{End-to-end latency.} The criterion is missed: the parse path's p95 from the end of
speech is 3{,}122~ms over 226 segments, against 2{,}500~ms, a shortfall of 622~ms (25\%). With the
endpointing wait included, the p95 is 3{,}602~ms. The per-stage rows locate the shortfall in two
stages. Speech recognition reaches 1{,}449~ms at p95 against its 1{,}200~ms allowance. Prefill
reaches 751~ms against 250~ms, because the deployed parser processes the whole prompt for every
command (Section~\ref{sec:runtime}), whereas the allowance assumes the fixed prefix is cached. The
decode, at 915~ms against 1{,}100~ms, meets its allowance. The selection rule of the
\emph{M\'emoire de Master} has a failure clause for a constraint that no configuration satisfies,
and under it the end-to-end budget is re-baselined to the measured 3{,}122~ms. This is a
re-baseline, not a pass: against the 2{,}500~ms requirement as specified, the selected
configuration misses by 622~ms, and the verdict of Section~\ref{sec:requirements-summary} remains
missed.

\paragraph{Per-stage attribution.} The criterion is met: each of the five parse-path stages of
Table~\ref{tab:stage-budget} is measured and reported at p50, p95 and p99. Three meet their
allowance: the endpointing wait (480~ms against 500~ms), the decode, and validation with the
state-machine check (under 1~ms against 50~ms). Two miss: speech recognition and prefill. Two
intervals inside the span have no budget row and are recorded beside the stages: the segment's
wait for the parse-path worker, 73~ms at p95, and the overhead of the requests to the
language-model server, 8~ms at p95. Two rows of Table~\ref{tab:stage-budget} have no verdict here.
The combined language-model allowance is the deployment constraint that the \emph{M\'emoire de
Master} applies to its model comparison, where the prefix is cached. The start-of-speech target is
defined for a 3~s utterance, while the golden-set utterances vary in length.

\paragraph{Keyword false rejects.} The author's 40 takes, scored at the operating threshold by the
keyword-spotter evaluation's own procedure, give one miss in 20 for the hold class, a false-reject
rate of 0.050 with an exact 95\% interval of [0.001, 0.249], and none in 20 for the abort class,
0.000 with [0.000, 0.168]. Under the declared rule the criterion is met only if the whole
interval lies below 0.10, and neither class's interval does. The criterion is therefore not
demonstrated. At 20 takes per class it could not have been: even with no miss, the upper bound
stays at 0.168. The one missed take peaked at a score of 0.9977 against the threshold of 0.999,
and it is the take missed on both replays on the board.

\paragraph{Not measured.} Two delays are outside every figure of this section: the latency of live
capture, since the audio was replayed, and the wireless hop from the board to the workstation,
since the bus ran over loopback. Section~\ref{sec:limitations} treats both.

\section{Keyword-spotter evaluation}
\label{sec:keyword-spotter-evaluation}
The keyword-spotter evaluation scores the two trained heads on held-out recordings, with the
operating threshold of 0.999 selected on the validation split before the test split was scored;
the test split was scored once for each of the heads' two fits (Section~\ref{sec:audio-chain}). It judges the keyword false-accept criterion and gives the
false-reject rates on synthetic speech that the author's recordings of
Section~\ref{sec:latency-experiment} complement. The recordings are streamed through the
spotter's front end in 80~ms frames, as the runtime streams them, and detections are counted with
the runtime's 1.0~s debounce. Every rate below therefore counts the same events the runtime would
publish.

\paragraph{Operating curve.} Figure~\ref{fig:keyword-curve} gives, for each class, the false-reject
rate on the test split against the false-accept rate on the ambient-speech stream, both classes
together, as the threshold sweeps the grid of Section~\ref{sec:audio-chain}, with the operating
point marked. The two keyword criteria are read from this one curve, as
Table~\ref{tab:nonfunctional-requirements} requires, and they are quoted together here. The
curve's upper end is flat in false accepts. A single false accept on the test stream survives
every threshold from 0.88 to the top of the grid, so no threshold on the grid would have removed
it; within that range, raising the threshold only adds false rejects. The operating point sits at
the top of the grid, where the heads' scores saturate and the threshold does little of the
classification.

\begin{figure}[htbp]
\centering
\figtodo{per-class curve from wake\_training.json, threshold 0.999 marked}
\caption[Keyword-spotter operating curve]{Operating curve of the keyword spotter on the held-out
test split. For each class, \texttt{swarm hold} and \texttt{swarm abort}, the false-reject rate on
the 300 synthetic test takes of that class is plotted against the false-accept rate per hour on
0.429~h of ambient speech, summed over both classes, as the common detection threshold sweeps the
selection grid from 0.05 to 0.999. The marker is the operating threshold of 0.999, selected on the
validation split before the test split was scored. Detections are counted in 80~ms frames with a
1.0~s debounce, as in the runtime.}
\label{fig:keyword-curve}
\end{figure}

\paragraph{False rejects on synthetic speech.} At the operating point, the hold class misses 0.027
of its 300 test takes and the abort class 0.007 of its 300. One hold take fired the abort class,
and no abort take fired the hold class. On the validation split, where the threshold was chosen,
the rates are 0.027 and 0.000 over 150 takes per class. Every positive of both splits is rendered
by the same three synthetic voices the heads were trained on, so these rates measure the heads on
voices they have heard. The criterion is judged on the author's voice
(Section~\ref{sec:latency-experiment}). The spotter's delay, from the last voiced sample of the
keyword to the first frame above the threshold, is 188~ms at p95 for the hold class and 131~ms for
the abort class on the test split. On the author's takes, measured the same way against the capture
tool's end of speech, it is 550~ms and 430~ms. The synthetic split therefore understates the delay on real
speech, which is where the reflex latency criterion is missed (Section~\ref{sec:latency-experiment}).

\paragraph{Keyword false accepts.} The test stream of ambient speech, 0.429~h long, produced one
false accept, on the abort head. That is 2.33 per hour, with a 95\% Poisson interval of [0.06,
12.98], against a budget of at most one per hour. The point estimate exceeds the budget, and the
interval contains it. Under the rule declared before scoring --- met only if the whole interval
lies inside the budget, missed only if it lies wholly outside --- the criterion is not
demonstrated. The stream is too short to decide it either way. The validation stream, 0.167~h,
produced no false accept, with an interval of [0.00, 22.13], and so decides nothing either. The
golden-set recordings, in which the author reads parse-path commands including phrases close to
the keywords, produced none in 0.304~h. A further stream of authored near-miss phrases, packed back
to back, produced 34 false accepts in 0.100~h, 31 of them on the hold head. That rate describes a
stress condition built from the nearest confusions available, not an operating rate, and no
criterion is judged on it. A false \texttt{abort} on the spoken word \emph{abort} is itself
fail-safe (Section~\ref{sec:membership-rule}).

\section{Acoustic-robustness experiment}
\label{sec:acoustic-robustness}
% Sources: results/exp3_pi.csv (the Pi run is of record); tab:end-to-end; results/
% exp3_pi_analysis.md (ANOVA/Tukey, Cochran's Q + McNemar, operating range, split rate, error
% analysis with one cause per failed utterance); results/exp3_pi_parity.csv.
% - Protocol: golden-set recordings with propeller-type noise mixed at falling SNR, through the
%   parse path from endpointing to the validator. A split utterance counts correct only if every
%   segment is.
% - Clean-audio recognition and recognition in noise: both missed, each with its measured cause.
%   Read exact match minus CRR as the cost of speech recognition.
% - Statistics: the pre-registered arcsine ANOVA with its two caveats (the transform on 0/1 data;
%   the same items repeated across SNR, hence Cochran's Q and McNemar beside it). The 10 dB vs
%   5 dB McNemar comparison is not significant, by a hair -- write it so (p and alpha from the
%   analysis file).
% - No tuning was attempted, and why: the golden set is the only real speech, so tuning on it
%   would fit the test set. Candidate remedies go to Ch 6.
% - Parity: the parser is bit-identical across the workstation and the Pi; speech recognition is
%   not, so accuracy is not hardware-independent and the target's numbers are reported; no verdict
%   or significant pair changed.
% - The stage-sum column of tab:end-to-end shows noise barely moves latency; end-to-end latency is
%   the latency experiment's, not this one's.
% - Safe-failure rate (the Memoire de Master's Safe failure criterion, same name): far below its
%   bound; how ADR-0006 re-baselined it. Ch 6 returns to it.
% - Thermal: the soft-limit flag was set during the 5 dB condition, checked, no measurable effect.
% - The WER on the author's voice: cite the Memoire de Master's speaker-sensitivity experiment.
\TODO{draft}

\section{Formation-control experiment}
\label{sec:formation-control}
% Sources: results/exp4.csv, results/exp4_formation.md; eval/exp4.py -> swarm/simulate.py
% (kinematic backend).
% - Protocol (fixed before measurement): circle, line, wedge; repeated 60 s trials, fixed seeds,
%   timestep and initial spread (simulate.py TrialSpec).
% - Collisions: zero observed across all trials, backed by a hard geometric separation clamp at
%   the integrator (Ch 1's locked phrasing). Beside it, the number of clamp interventions -- the
%   informative figure; zero follows from the clamp's convergence, since collisions are counted
%   after the clamp at a distance below the clamp distance.
% - Formation accuracy: saturated, and this was declared before, not discovered after;
%   convergence time is the discriminating metric; the ANOVA on it.
% - Scope: kinematic simulator without rotor dynamics; the clamp's scope is Ch 6's limitation.
\TODO{draft}

\section{Summary against the requirements}
\label{sec:requirements-summary}
% Author's decision (24 Sep): one closing table, every criterion of Ch 1's two requirement tables
% with its target, measurement, verdict and section. prd.md 3.1 item 5 amended to list it.
% - % Generated by eval/tables.py -- do not edit; edit the renderer and re-run.
\begin{table}[htbp]
  \centering
  \footnotesize
  \setlength{\tabcolsep}{4pt}
  \caption[Summary against the requirements]{Evaluation summary of all functional and non-functional requirements in order, stating the answering measurement, reported section, and verdict.}
  \label{tab:requirements-summary}
  \begin{tabular}{@{}>{\raggedright\arraybackslash}p{3.5cm}>{\raggedright\arraybackslash}p{3.6cm}>{\raggedright\arraybackslash}p{4.6cm}>{\raggedright\arraybackslash}p{2.5cm}r@{}}
    \toprule
    Criterion & Target & Measured & Verdict & Section \\
    \midrule
    Offline speech recognition & On the device, no network & No run with networking disabled & Not yet run & \ref{sec:demonstration-protocol} \\
    Reflex path & Preempts the parse path & 78/78 decodes cancelled; ordering rule unit-tested; 0 cross-triggers & Met & \ref{sec:latency-experiment} \\
    Swarm controller & Formation control, five drones & 150/150 trials converged & Met & \ref{sec:formation-control} \\
    Simulation backends & Hover smoke test on both & Passes on both, same gains & Met & \ref{sec:simulation-backends} \\
    Live demonstration & Live microphone, end to end & -- & Planned for the defence & \ref{sec:demonstration-protocol} \\
    Flight state machine & Every legality cell tested & Every cell passes & Met & \ref{sec:state-machine} \\
    Reflex latency & p95 $\leq$ 150 ms & 545 ms idle, 547 ms loaded & Missed & \ref{sec:latency-experiment} \\
    End-to-end latency & p95 $\leq$ 2,500 ms & 3,122 ms & Missed & \ref{sec:latency-experiment} \\
    Per-stage attribution & Every stage, p50/p95/p99 & All 5 stages & Met & \ref{sec:latency-experiment} \\
    Preemption recovery & p95 $\leq$ 300 ms & 1,195 ms & Missed & \ref{sec:latency-experiment} \\
    Keyword false accepts & $\leq$ 1 per hour & 2.33/h [0.06, 12.98] & Not demonstrated & \ref{sec:keyword-spotter-evaluation} \\
    Keyword false rejects & $\leq$ 0.10 per class & hold 0.05 [0.00, 0.25]; abort 0.00 [0.00, 0.17] & Not demonstrated & \ref{sec:latency-experiment} \\
    Clean-audio recognition & CRR $\geq$ 0.80 & 0.690 & Missed & \ref{sec:acoustic-robustness} \\
    Recognition in noise & CRR $\geq$ 0.65 at 10 dB & 0.590 & Missed & \ref{sec:acoustic-robustness} \\
    Collisions & Zero observed & 0 in 150 trials & Met & \ref{sec:formation-control} \\
    Formation accuracy & $\geq$ 0.85 & 1.000 in every trial & Met & \ref{sec:formation-control} \\
    Offline operation & No network dependency & No run with networking disabled & Not yet run & \ref{sec:demonstration-protocol} \\
    \bottomrule
  \end{tabular}
  \par\medskip
  \begin{minipage}{\textwidth}\footnotesize\raggedright Collisions: zero observed across all 150 trials, backed by a hard geometric separation clamp at the integrator, which made 432 interventions (vehicle pairs, summed over timesteps, that the potential field alone failed to keep apart). Formation accuracy is saturated by the protocol; convergence time is the measure that discriminates (Section~\ref{sec:formation-control}).\end{minipage}
\end{table}
 (\label{tab:requirements-summary}, from
%   eval/tables.py requirements_summary()). Never retype it; a wording change goes in the renderer.
%   eval/test_tables.py fails if its rows stop matching Ch 1's criteria, in order.
% - Prose, short: how many criteria are met, missed, not demonstrated, not yet run -- count them
%   from the table, do not type the counts from this note. Name the misses in one sentence each
%   with their located cause, pointing back to the section; no new figure, no new argument.
% - The verdict vocabulary is defined in the caption: point estimate for Met/Missed; the declared
%   interval rule for the two keyword criteria ("not demonstrated"). Use the same words in prose.
% - This section is what Ch 6 refers to for every verdict; it does not interpret (Ch 6 does).
\TODO{draft}
